\documentclass[11pt]{article}
\usepackage[a4paper,margin=1in]{geometry}
\usepackage{fontspec}
\usepackage[boldfont=false]{xeCJK}
\setCJKmainfont{FandolSong-Regular.otf}
\usepackage{amsmath}
\usepackage{amssymb}
\usepackage{array}
\usepackage{booktabs}
\usepackage{graphicx}
\usepackage{adjustbox}
\usepackage{enumitem}
\setlist[itemize]{topsep=2pt,partopsep=0pt,parsep=0pt,itemsep=2pt,leftmargin=1.4em}
\setlist[enumerate]{topsep=2pt,partopsep=0pt,parsep=0pt,itemsep=2pt,leftmargin=1.6em}
\usepackage{parskip}
\usepackage{fvextra}
\fvset{breaklines=true,breakanywhere=true,fontsize=\small}
\setlength{\parindent}{0pt}

\begin{document}

\begin{center}
  {\LARGE\bfseries Cell structure}\\[0.35em]
  {\large A Level Biology}
\end{center}
\vspace{0.6em}

\section*{How we study cells}

Cells 细胞 are very small, so you cannot see them with your eyes alone. You use a \textbf{microscope} 显微镜 to make a bigger picture of them. The first kind you meet is the \textbf{light microscope} 光学显微镜. It shines light through a thin \textbf{specimen} 标本 (the material you look at) and uses glass lenses to enlarge the view.

\subsection*{Making a slide and drawing what you see}

To look at living material, you make a \textbf{temporary preparation} 临时装片. You put a small, thin piece of material on a glass \textbf{slide} 载玻片, add a drop of \textbf{stain} 染色剂 (a coloured liquid that makes parts easier to see), then lower a thin \textbf{cover slip} 盖玻片 on top to flatten it and keep out air.

When you draw cells from a slide or a photograph, follow simple rules:

\begin{itemize}
\item use a sharp pencil and clear, single lines (no shading).

\item draw only what you can really see, with the parts in the correct sizes.

\item label the parts with straight lines that do not cross.

\end{itemize}

\subsection*{Magnification and actual size}

\textbf{Magnification} 放大倍数 tells you how many times bigger the image is than the real object. It has no unit. You find it with one equation:

\[
\text{magnification} = \frac{\text{size of image}}{\text{actual size of object}}
\]

You can rearrange this to find any one value from the other two:

\[
\text{actual size} = \frac{\text{size of image}}{\text{magnification}}
\]

The top and the bottom of the fraction \textbf{must use the same unit}. Cells are tiny, so you work in small units:

\begin{itemize}
\item $1\ \text{mm} = 1000\ \text{micrometre}$ 微米 (µm)

\item $1\ \text{µm} = 1000\ \text{nanometre}$ 纳米 (nm)

\item so $1\ \text{mm} = 1\,000\,000\ \text{nm}$.

\end{itemize}

\textbf{Worked example.} In a photomicrograph at magnification $5000$, a \textbf{chloroplast} 叶绿体 measures $25\ \text{mm}$ across. Its actual size is

\[
\frac{25\ \text{mm}}{5000} = 0.005\ \text{mm} = 5\ \text{µm}.
\]

The same equation works for drawings, photomicrographs 显微照片 (photos taken through a light microscope) and \textbf{electron micrographs} 电子显微照片 (the most detailed photos, explained below). Always convert to the same unit first, then divide.

\subsection*{Eyepiece graticule and stage micrometer}

To measure a real cell under the microscope, you use an \textbf{eyepiece graticule} 目镜测微尺 — a tiny scale inside the eyepiece. Its divisions have no fixed size, so first you must \textbf{calibrate} 校准 them (work out what one division is worth).

You calibrate using a \textbf{stage micrometer} 载物台测微尺 — a special slide with an accurate scale on it (often $1\ \text{mm}$ split into $100$ parts, so each part is $10\ \text{µm}$). You line up the two scales, count how many graticule divisions fit a known length, and divide. Once calibrated, you can swap in your specimen and measure it with the graticule.

\subsection*{Resolution and magnification}

These two words are easy to mix up. The examiner gives marks for the difference.

\begin{itemize}
\item \textbf{Magnification} is how many times bigger the image is than the object.

\item \textbf{Resolution} 分辨率 is the smallest distance between two points that still lets you see them as \textbf{two} separate points.

\end{itemize}

Making an image bigger does not always show more detail. Past a certain point you just get a bigger, blurry image. Resolution sets the real limit on detail.

A light microscope has lower resolution because light has a fairly long \textbf{wavelength} 波长. An \textbf{electron microscope} 电子显微镜 uses beams of \textbf{electrons} 电子 instead of light. Electrons have a much shorter wavelength, so the resolution is far higher and you can see very small structures inside the cell. There are two kinds: \textbf{scanning} 扫描 (shows the surface in 3D) and \textbf{transmission} 透射 (passes electrons through a thin slice to show inside detail).

\section*{Eukaryotic cells and their organelles}

Plant and animal cells are \textbf{eukaryotic cells} 真核细胞: their DNA is kept inside a \textbf{nucleus} 细胞核. Inside the cell are many small parts called \textbf{organelles} 细胞器, each with its own job. The jelly-like fluid around them is the \textbf{cytoplasm} 细胞质.

In a photomicrograph or electron micrograph you identify organelles by their shape, size and position; in a drawing you show their outlines and label them.

\begin{center}
\adjustbox{max width=\linewidth}{%
\begin{tabular}{lll}
\toprule
\textbf{Organelle} & \textbf{Structure} & \textbf{Function} \\
\midrule
\textbf{cell surface membrane} 细胞膜 & thin layer around the cell & controls what enters and leaves the cell \\
nucleus & large, surrounded by a \textbf{nuclear envelope} 核膜 (a double membrane with holes); contains a \textbf{nucleolus} 核仁 & holds the DNA; controls the cell; the nucleolus makes ribosomes \\
\textbf{rough endoplasmic reticulum} 粗面内质网 (rough ER) & sheets of membrane with ribosomes on the surface & makes and transports \textbf{proteins} 蛋白质 (for example \textbf{antibodies} 抗体) \\
\textbf{smooth endoplasmic reticulum} 滑面内质网 (smooth ER) & sheets of membrane, no ribosomes & makes \textbf{lipids} 脂质 \\
\textbf{Golgi body} 高尔基体 & stack of flat membrane sacs & changes and packs proteins and lipids into \textbf{vesicles} 囊泡 for \textbf{secretion} 分泌 \\
\textbf{mitochondria} 线粒体 & oval, with a folded inner membrane; has small circular DNA & site of \textbf{respiration} 呼吸作用 — releases energy 能量 as ATP \\
\textbf{ribosomes} 核糖体 & very small; $80\text{S}$ in the cytoplasm, $70\text{S}$ in chloroplasts and mitochondria & join \textbf{amino acids} 氨基酸 to \textbf{synthesise} 合成 proteins \\
\textbf{lysosomes} 溶酶体 & small sacs of \textbf{enzymes} 酶 & break down old organelles and waste \\
\textbf{centrioles} 中心粒 and \textbf{microtubules} 微管 & small tubes made of protein & help move chromosomes and form the cell's "skeleton" \\
\textbf{cilia} 纤毛 & tiny hairs on the cell surface that beat & move fluid or move the cell \\
\textbf{microvilli} 微绒毛 & tiny folds of the cell surface membrane & increase surface area for \textbf{absorption} 吸收 \\
\textbf{chloroplasts} (plants) & green, with stacked membranes; has small circular DNA & site of \textbf{photosynthesis} 光合作用 \\
\textbf{cell wall} 细胞壁 (plants) & strong outer layer of \textbf{cellulose} 纤维素 & supports and protects the cell; stops it bursting \\
\textbf{plasmodesmata} 胞间连丝 (plants) & tiny channels through the cell walls & link the cytoplasm of neighbouring cells \\
large permanent \textbf{vacuole} 液泡 (plants) & big sac of watery fluid, with a membrane called the \textbf{tonoplast} 液泡膜 & stores water and keeps the cell firm \\
\bottomrule
\end{tabular}}
\end{center}

Cells use \textbf{ATP} made in respiration as their energy supply for every job that needs energy, such as making proteins, moving things and dividing.

\subsection*{Comparing plant and animal cells}

\begin{center}
\adjustbox{max width=\linewidth}{%
\begin{tabular}{lll}
\toprule
\textbf{Feature} & \textbf{Plant cell} & \textbf{Animal cell} \\
\midrule
cell wall & present (cellulose) & absent \\
chloroplasts & present & absent \\
large permanent vacuole & present & absent (only small, temporary ones) \\
centrioles & absent in most & present \\
shape & fixed and regular & rounder and more flexible \\
\bottomrule
\end{tabular}}
\end{center}

Both have a cell surface membrane, cytoplasm, a nucleus, mitochondria, ribosomes, ER and a Golgi body.

\section*{Prokaryotic cells (bacteria)}

A \textbf{prokaryotic cell} 原核细胞, such as a \textbf{bacterium} 细菌, is much smaller and simpler than a eukaryotic cell. Its key features are:

\begin{itemize}
\item \textbf{unicellular} 单细胞 — it is a single cell.

\item generally $1$–$5\ \text{µm}$ across.

\item a cell wall made of \textbf{peptidoglycan} 肽聚糖 (not cellulose).

\item circular DNA lying free in the cytoplasm — there is \textbf{no nucleus}.

\item $70\text{S}$ ribosomes (smaller than the $80\text{S}$ ones in the cytoplasm of eukaryotes).

\item \textbf{no organelles surrounded by a double membrane} — so no nucleus, no mitochondria and no chloroplasts.

\end{itemize}

\subsection*{Comparing prokaryotic and eukaryotic cells}

\begin{center}
\adjustbox{max width=\linewidth}{%
\begin{tabular}{lll}
\toprule
\textbf{Feature} & \textbf{Prokaryotic cell} & \textbf{Eukaryotic cell} \\
\midrule
size & about $1$–$5\ \text{µm}$ & about $10$–$100\ \text{µm}$ \\
DNA & circular, free in cytoplasm & linear, inside a nucleus \\
nucleus & none & present \\
double-membrane organelles & none & mitochondria (and chloroplasts in plants) \\
ribosomes & $70\text{S}$ & $80\text{S}$ (with $70\text{S}$ inside mitochondria and chloroplasts) \\
cell wall & peptidoglycan & cellulose (plants) or none (animals) \\
\bottomrule
\end{tabular}}
\end{center}

\section*{Viruses}

All \textbf{viruses} 病毒 are \textbf{non-cellular} 非细胞 — they are not made of cells at all. Each virus is built from just two or three parts:

\begin{itemize}
\item a core of \textbf{nucleic acid} 核酸, which is \textbf{either DNA or RNA} (never both).

\item a protein coat around the core called a \textbf{capsid} 衣壳.

\item in some viruses, an outer \textbf{envelope} 包膜 made of \textbf{phospholipids} 磷脂.

\end{itemize}

A virus has no cytoplasm, no organelles and no ribosomes. It cannot respire or make its own proteins. It can only copy itself inside a living \textbf{host} 宿主 cell, so it sits at the edge of what we call "living".


\end{document}
